%==============================================================================
% Sjabloon poster graduaatsproject
%==============================================================================
% Gebaseerd op de documentclass `a0poster' van Gerlinde Kettl en Matthias Weiser
% Aangepast voor HOGENT door Jens Buysse en Bert Van Vreckem
%
% Je gegevens staan in ../projectgegevens.tex.
%
% De poster is een affiche van A0-formaat: hij wordt van op twee meter afstand
% bekeken, op de verdedigingsdag en op infodagen. Schrijf dus kort, gebruik
% beeld, en zorg dat iemand in dertig seconden begrijpt wat je gemaakt hebt.

\documentclass[a0,portrait]{hogent-poster}

%%=============================================================================
%% Projectgegevens 
%%=============================================================================

%% ---------- Jij ------------------------------------------------------------
\newcommand{\studentnaam}{Sarah Turner Burkitt}
\newcommand{\studentemail}{sarah.turnerburkitt@student.hogent.be}

%% Wordt gebruikt als bestandsnaam van je PDF's, bijvoorbeeld
%% Familienaam-Voornaam-rapport.pdf. Geen spaties of leestekens gebruiken.
\newcommand{\bestandsnaam}{TurnerBurkitt-Sarah}

%% ---------- Je project -----------------------------------------------------
\newcommand{\projecttitel}{Gerichte configuratie van automatisch gegenereerde output}
\newcommand{\projectondertitel}{Uitbreiding van document- en consultatietemplates binnen Vesalius.ai}
%% Een of twee zinnen. Wordt gebruikt op de poster en in het voorstel.
\newcommand{\projectpitch}{%
    Dit graduaatsproject breidt het Vesalius-platform uit met gerichtere configuratie van automatisch gegenereerde output: 
    documenttemplates kunnen aan specifieke artsen worden gekoppeld en Scribe-consultatemplates kunnen worden afgestemd op patiënt- 
    of artsgerichte output.%
}

%% ---------- Waar je het doet ------------------------------------------------
\newcommand{\bedrijfsnaam}{Vesalius.ai}

%% Je mentor op de werkvloer (het bedrijf)
\newcommand{\mentornaam}{Bart D'Hoore}
\newcommand{\mentoremail}{Bart.D'Hoore@endare.be}

%% Je begeleider van de opleiding (HOGENT)
\newcommand{\begeleidernaam}{Luc Vervoort}
\newcommand{\begeleideremail}{luc.vervoort@hogent.be}

%% ---------- Je repository ---------------------------------------------------
\newcommand{\repolink}{https://github.com/STB95/Graduaatsproject_TurnerBurkittSarah_26-27}
%% op je poster: hij is een beoordeeld onderdeel van je oplevering.
\newcommand{\repolink}{https://les.lucvervoort.org/sarahturnerburkitt/graduaatsproject}

%% ---------- Vast voor dit opleidingsonderdeel -------------------------------
%% Deze twee hoef je normaal niet aan te passen.
\newcommand{\opleidingnaam}{Graduaat in het Programmeren}
\newcommand{\academiejaarwaarde}{2026-2027}

%% Zoekwoorden voor het voorstel en de poster (3 tot 5)
\newcommand{\projectkeywords}{Vesalius.ai, documenttemplates, Scribe, automatische output, Angular/Laravel, WYSIWYG e-maileditor}


% Info over de opleiding
\course{Graduaatsproject}
\studyprogramme{\opleidingnaam}
\academicyear{\academiejaarwaarde}
\institution{Hogeschool Gent, Valentin Vaerwyckweg 1, 9000 Gent}

% Info over het project
\title{\projecttitel}
\subtitle{\projectondertitel}
\author{\studentnaam}
\email{\studentemail}
\supervisor{\begeleidernaam}
\cosupervisor{\mentornaam{} (\bedrijfsnaam)}

% Deze informatie komt onderaan de samenvatting te staan
\specialisation{Programmeren}
\keywords{\projectkeywords}
\projectrepo{\repolink}

\begin{document}

\maketitle

\begin{abstract}
  % Drie tot vijf zinnen: voor welk bedrijf, welk probleem, wat heb je gemaakt
  % en wat levert het op. Dit is het enige stuk lopende tekst op je poster dat
  % iemand echt van begin tot eind leest.
  Vervang deze tekst door de samenvatting van je project.
\end{abstract}

\begin{multicols}{2}

\section{Probleem}

% Wat liep er mis of ontbrak er, en voor wie? Twee tot vier korte zinnen, of
% drie bullets. Geen achtergrondverhaal.

\section{Aanpak en technologie}

% Hoe heb je het aangepakt en waarmee? Noem de technologie bij naam en zeg in
% één zin waarom het die geworden is. Een klein architectuurschema doet hier
% meer werk dan een paragraaf tekst.

\begin{center}
  \captionsetup{type=figure}
  \includegraphics[width=1.0\linewidth]{voorbeeldafbeelding}
  \captionof{figure}{Vervang deze afbeelding door je architectuurschema of
    een schermafbeelding van je oplossing. Op een poster is beeld belangrijker
    dan tekst: reken op minstens twee afbeeldingen.}
\end{center}

\section{Resultaat}

% Wat is er opgeleverd, en werkt het? Toon het liefst met een schermafbeelding
% of een cijfer: "verwerkt 1200 orders per uur", "bespaart het supportteam
% ongeveer een uur per dag". Concreet verslaat algemeen.

\section{Wat ik eruit meeneem}

% Twee tot drie zinnen: wat heb je geleerd, technisch en over jezelf als
% beginnende professional? Dit is het menselijke stuk van je poster, en het is
% wat bezoekers op een infodag onthouden.

\end{multicols}
\end{document}
